% TOP_PROBLEM: {"id": 30006713, "title": "Chowla conjecture on correlations of the Liouville function", "category_id": 1, "category": {"id": 1, "name": "number_theory", "display_name": "Number Theory", "description": "Properties of integers, prime numbers, Diophantine equations.", "slug": "number-theory", "order_index": 1, "created_at": "2026-07-31T15:26:25.670Z"}, "status": "open"}
% FIRST_POSTED: 2026-09-25
% ENTRY_KIND: statement_only
% !TeX program = lualatex
% Definition and sources only; this file is not an LLM solution attempt.
\documentclass[11pt]{article}
\usepackage[margin=25mm]{geometry}
\usepackage{fontspec}
\setmainfont{Latin Modern Roman}
\usepackage{amsmath,amssymb,mathtools}
\usepackage{unicode-math}
\setmathfont{Latin Modern Math}
\usepackage{xurl,hyperref}
\hypersetup{colorlinks=true,urlcolor=blue}
\setcounter{secnumdepth}{0}
\setlength{\parindent}{0pt}
\setlength{\parskip}{5pt}
\setlength{\emergencystretch}{3em}
\begin{document}
\section*{\texorpdfstring{Chowla conjecture on correlations of the Liouville function}{Chowla conjecture on correlations of the Liouville function}}\label{30006713}
\subsection{Definitions and mathematical statement}
For $n=\prod_p p^{v_p(n)}$, put $\Omega(n)=\sum_pv_p(n)$ and $\lambda(n)=(-1)^{\Omega(n)}$. For the positive-slope linear forms $L_i(n)=a_in+b_i$ specified in the record, require pairwise nonproportionality $a_ib_j-a_jb_i\ne0$. The assertion is
\[
 \frac1x\sum_{1\le n\le x}\prod_{i=1}^k\lambda(a_in+b_i)\longrightarrow0.
\]
This is an ordinary, unweighted Cesàro average. Replacing it by a logarithmically weighted average changes the claim. The record's distinct nonnegative offsets and positive natural slopes are retained.
\par\smallskip\textit{Formulation and concept references:} \href{https://arxiv.org/abs/1605.04628}{[S1]}.

\subsection{Short English statement}
Do the parities of prime-factor counts decorrelate along every allowed finite family of distinct linear patterns?
\subsection{Sources}
\begin{itemize}
\item[S1]Terence Tao, Equivalence of the logarithmically averaged Chowla and Sarnak conjectures, abstract (2016).
\url{https://arxiv.org/abs/1605.04628}.
\textit{Formulation source.}
\item[S2]The Chowla conjecture and Landau–Siegel zeroes.
\url{https://www.cambridge.org/core/journals/mathematical-proceedings-of-the-cambridge-philosophical-society/article/chowla-conjecture-and-landausiegel-zeroes/515F3378450DED201A21E26BF801CEEB}.
\textit{Further reference; not a proof endorsement.}
\item[S3]From Bateman-Horn to Chowla.
\url{https://www.preprints.org/manuscript/202601.0277}.
\textit{Further reference; not a proof endorsement.}
\item[S4]Two averaged dynamical generalizations of Chowla's conjecture.
\url{https://arxiv.org/abs/2608.16108}.
\textit{Further reference; not a proof endorsement.}
\end{itemize}
\end{document}
